\lecture{Лекция 1}{}

\section*{Введение. Мотивирующая задача}
\addcontentsline{toc}{section}{Введение. Мотивирующая задача}

\subsection*{Можно ли восстановить функцию по её производным в точке?}

\textbf{Вопрос.} Пусть функция $f$ бесконечно дифференцируема в окрестности точки $x_0$ и
\[
  f(x_0) = 0, \qquad f^{(k)}(x_0) = 0 \quad \forall k \in \N .
\]
Чему равно $f(x_0 + \Delta x)$?

\textbf{Наивный ответ} — «нуль». Он получается из формулы Тейлора:
\[
  f(x_0 + \Delta x) = \sum_{k=0}^{\infty} \frac{f^{(k)}(x_0)}{k!}\,(\Delta x)^k = 0 .
\]

\textbf{Правильный ответ:} \emph{вообще говоря, ничего сказать нельзя}. Ряд Тейлора
действительно тождественно равен нулю, но гладкая функция \textbf{не обязана совпадать
со своим рядом Тейлора}. Совпадение гарантировано только для \emph{аналитических}
функций, а класс $C^\infty$ гораздо шире.

\begin{idea}
У аналитических функций (многочленов, $e^x$, $\sin x$ и т.\,п.) все производные в одной
точке определяют функцию в окрестности этой точки. Для функций из $C^\infty$ это неверно:
бесконечная гладкость — свойство гораздо более слабое, чем аналитичность. Именно эта
«гибкость» гладких функций и будет нам нужна.
\end{idea}

\subsection*{Контрпример: «функция-шапочка»}

\[
  f(x) =
  \begin{cases}
    e^{\frac{1}{x^2 - 1}}, & -1 < x < 1, \\[4pt]
    0, & |x| \ge 1 .
  \end{cases}
\]

\begin{warn}
На доске во второй строке было написано $x = \pm 1$. Из дальнейшего ясно, что имелось
в виду $|x| \ge 1$: функция продолжается нулём на всю прямую.
\end{warn}

Ту же функцию часто записывают в виде $f(x) = e^{-\frac{1}{1 - x^2}}$ при $|x| < 1$.

\textbf{Свойства:}
\begin{enumerate}
  \item \textbf{$f(x) > 0$ при $|x| < 1$}, например $f(0) = e^{-1}$.
  \item \textbf{$f$ непрерывна в точках $\pm 1$.} При $x \to \pm 1$ изнутри
        $\frac{1}{x^2-1} \to -\infty$, поэтому $f(x) \to 0$.
  \item \textbf{$f \in C^\infty(\R)$, причём $f^{(k)}(\pm 1) = 0$ для всех $k$.}
        По индукции производные внутри интервала имеют вид
        \[
          f^{(k)}(x) = \frac{P_k(x)}{(x^2-1)^{2k}}\; e^{\frac{1}{x^2-1}},
        \]
        где $P_k$ — многочлен. После замены $t = \frac{1}{1-x^2} \to +\infty$ получается
        выражение вида «многочлен от $t$, умноженный на $e^{-t}$», и оно стремится к $0$.
        Значит, все односторонние производные в точках $\pm 1$ нулевые и совпадают
        с производными нулевой функции снаружи.
\end{enumerate}

\textbf{Вывод.} В точке $x_0 = 1$ функция и все её производные равны нулю,
но $f(x_0 + \Delta x) \neq 0$ при $-2 < \Delta x < 0$.

\begin{idea}
\textbf{Зачем это в курсе.} Это пример \textbf{гладкой функции с компактным носителем}
(класс $C_c^\infty$, он же $C_0^\infty$ или $\mathcal{D}$). Аналитическая функция так себя
вести не может: если она равна нулю на интервале, то она равна нулю всюду. Функции из
$C_0^\infty$ — главный рабочий инструмент курса. Из «шапочки» строятся ядра усреднения
(глава~1), через $C_0^\infty$ определяются обобщённые производные.
\end{idea}
